\subsection{\texorpdfstring{第~\ref{sec:special-relativity}~节}{狭义相对论}}\label{subsec:solutions-special-relativity}

\begin{description}[font=\normalfont,leftmargin=2em,labelwidth=1.5em]
\RGsolution{ex:signature-of-nondegenerate-symmetric-form}{signature-of-nondegenerate-symmetric-form}
因为 $\langle\cdot,\cdot\rangle$ 非退化, 存在 $v,w\in V$ 使 $\langle v,w\rangle\ne0$. 此外, 存在 $v\in V$ 使 $\langle v,v\rangle\ne0$. 实际上, 即便 $\langle v,v\rangle=\langle w,w\rangle=0$, 也可取 $\widetilde v:=v+w$, 此时
\[
  \langle v+w,v+w\rangle=2\langle v,w\rangle\ne0.
\]
现在对 $V$ 的维数 $n$ 归纳, 证明存在标准正交基. 若 $n=1$, 取 $w\in V$ 使 $\langle w,w\rangle\ne0$, 并定义
\[
  v_1:=\frac{w}{|w|},
\]
其中 $|w|:=|\langle w,w\rangle|^{1/2}$. 显然 $\langle v_1,v_1\rangle=\pm1$, 且 $\{v_1\}$ 是所需的标准正交基.

若 $n>1$, 又取 $w\in V$ 使 $\langle w,w\rangle\ne0$, 并令
\[
  v_1:=\frac{w}{|w|}.
\]
若 $W$ 是 $v_1$ 张成的空间在 $V$ 中的正交补, 则 $\dim W=n-1$, $v_1\notin W$, 且 $\langle\cdot,\cdot\rangle$ 在 $W$ 上的限制是对称且非退化的. 应用归纳假设得到 $W$ 的一组基 $\{v_2,\ldots,v_n\}$, 使得 $i\ne j$ 时 $\langle v_i,v_j\rangle=0$, 且当 $i=2,\ldots,n$ 时 $\langle v_i,v_i\rangle=\pm1$. 因此 $\{v_1,\ldots,v_n\}$ 是 $V$ 的所需的基.

为说明 $\langle\cdot,\cdot\rangle$ 的符号差不依赖于标准正交基的选取, 注意它可以不依赖基地定义为: $\langle\cdot,\cdot\rangle$ 正定的 $V$ 的一个极大子空间的维数, 减去 $\langle\cdot,\cdot\rangle$ 负定的 $V$ 的一个极大子空间的维数.

\RGsolution{ex:minkowski-orthogonality-and-time-orientation}{minkowski-orthogonality-and-time-orientation}
固定惯性坐标 $(x^0,x^1,x^2,x^3)$. 则
\[
  v=\sum_{i=0}^3v^i\frac{\partial}{\partial x^i}
  \quad\text{且}\quad w=\sum_{i=0}^3w^i\frac{\partial}{\partial x^i}.
\]
\begin{enumerate}[label=(\alph*),ref=\alph*,leftmargin=*]
\item
\begin{enumerate}[label=(\roman*),ref=\roman*,leftmargin=*]
\item\label{sol:inner-product-with-future-timelike-vector} 因为 $v$ 类时且指向未来, 有
\[
  \langle v,v\rangle=-(v^0)^2+\sum_{i=1}^3(v^i)^2<0\quad\text{且}\quad v^0>0.
\]
类似地, 因为 $w$ 类时或零模且指向未来, 有
\[
  \langle w,w\rangle=-(w^0)^2+\sum_{i=1}^3(w^i)^2\le0\quad\text{且}\quad w^0>0.
\]
于是由 Cauchy–Schwarz 不等式,
\begin{align*}
  \langle v,w\rangle&=-v^0w^0+\sum_{i=1}^3v^iw^i\\
  &\le -v^0w^0+\left(\sum_{i=1}^3(v^i)^2\right)^{1/2}\left(\sum_{i=1}^3(w^i)^2\right)^{1/2}\\
  &<-v^0w^0+|v^0||w^0|=0.
\end{align*}
\item 由于
\[
  \langle v+w,v+w\rangle=\langle v,v\rangle+2\langle v,w\rangle+\langle w,w\rangle
\]
且 $\langle v,v\rangle<0$, $\langle w,w\rangle\le0$, $\langle v,w\rangle<0$ [由 (\ref{sol:inner-product-with-future-timelike-vector})], 所以 $\langle v+w,v+w\rangle<0$. 此外,
\[
  (v+w)^0=v^0+w^0>0,
\]
因此 $v+w$ 也指向未来.
\item 由 (\ref{sol:inner-product-with-future-timelike-vector}), 若 $\langle v,w\rangle=0$, 则 $w$ 不可能类时或零模 (零向量除外).
\end{enumerate}
\item
\begin{enumerate}[label=(\roman*),ref=\roman*,leftmargin=*]
\item\label{sol:inner-product-with-future-null-vector} 因为 $v$ 零模且指向未来, 有 $(v^0)^2=\sum_{i=1}^3(v^i)^2$ 且 $v^0>0$. 此外, 因为 $w$ 类时或零模且指向未来, 有 $\bigl(\sum_{i=1}^3(w^i)^2\bigr)^{1/2}\le w^0$ 且 $w^0>0$. 于是由 Cauchy–Schwarz 不等式,
\begin{align*}
  \langle v,w\rangle&=-v^0w^0+\sum_{i=1}^3v^iw^i\\
  &\le -v^0w^0+\left(\sum_{i=1}^3(v^i)^2\right)^{1/2}\left(\sum_{i=1}^3(w^i)^2\right)^{1/2}\\
  &\le-v^0w^0+|v^0||w^0|=0.
\end{align*}
此外, 等号成立当且仅当
\[
  \sum_{i=1}^3v^iw^i=\left(\sum_{i=1}^3(v^i)^2\right)^{1/2}
       \left(\sum_{i=1}^3(w^i)^2\right)^{1/2}
\]
且
\[
  \left(\sum_{i=1}^3(w^i)^2\right)^{1/2}=w^0.
\]
因此, 等号成立当且仅当 $w^i=\lambda v^i$, 其中 $\lambda>0$, 且
\[
  w^0=\lambda\left(\sum_{i=1}^3(v^i)^2\right)^{1/2}=\lambda v^0.
\]
\item 因为
\[
  \langle v+w,v+w\rangle=2\langle v,w\rangle+\langle w,w\rangle
\]
且 $\langle w,w\rangle\le0$ 及 $\langle v,w\rangle\le0$ [由 (\ref{sol:inner-product-with-future-null-vector})], 所以 $\langle v+w,v+w\rangle\le0$. 此外, 等号成立当且仅当 $\langle v,w\rangle=\langle w,w\rangle=0$, 即当且仅当 $w=\lambda v$, 其中 $\lambda>0$. 在所有情形下, $v+w$ 都显然指向未来.
\item 由 (\ref{sol:inner-product-with-future-null-vector}), 若 $\langle v,w\rangle=0$, 则 $w$ 不可能类时; 若 $w$ 零模, 它就是 $v$ 的倍数.
\end{enumerate}
\item 因为 $v$ 类空, 必有 $A:=\sum_{i=1}^3(v^i)^2>(v^0)^2\ge0$. 则向量 $w:=(A,v^0v^1,v^0v^2,v^0v^3)$ 类时, 且 $\langle v,w\rangle=0$.

此外, $v^1,v^2,v^3$ 中至少有一个不为零. 不妨设 $v^1\ne0$, 则向量 $\widetilde w=(0,-v^2,v^1,0)$ 类空, 且满足 $\langle v,\widetilde w\rangle=0$. 此外, $\langle w,\widetilde w\rangle=0$, 所以
\[
  \langle aw+b\widetilde w,aw+b\widetilde w\rangle
  =a^2\langle w,w\rangle+b^2\langle\widetilde w,\widetilde w\rangle.
\]
因为 $w$ 类时而 $\widetilde w$ 类空, 可取非零的 $a$ 和 $b$ 使这个量为零.
\end{enumerate}

\RGsolution{ex:twin-paradox}{twin-paradox}
\begin{enumerate}[label=(\alph*),ref=\alph*,leftmargin=*]
\item 取年和光年分别作为时间与长度单位. 在地球参考系中, 地球的运动是类时直线 $x=0$, 而行星 X 的运动是类时直线 $x=8$. 若选取 $t=0$ 使 Bob 的出发为事件 $(0,0)$, 则 Bob 到达行星 X 为事件 $(10,8)$, 双胞胎重逢为事件 $(20,0)$. Bob 的运动是连接这些事件的折线, 因此 Bob 在出发与返回之间测得的时间为
\begin{align*}
  |(10,8)-(0,0)|+|(20,0)-(10,8)|
  &=|(10,8)|+|(10,-8)|\\
  &=\sqrt{10^2-8^2}+\sqrt{10^2-(-8)^2}=6+6=12.
\end{align*}
因此, Bob 再次见到 Alice 时年龄为 $20+12=32$ 岁.
\item 虽然 Bob 的确可以说在他的参考系中是 Alice 在运动, 但他的参考系不是惯性参考系, 因为他须在事件 $(10,8)$ 处加速以反转速度. 因此不能在 Bob 的参考系中使用 Minkowski 几何.
\item 在事件 $(10,8)$ 处, Bob 正接收到 $t=2$ 时离开地球的光. 因此, 在到达行星 X 的 6 年中, Bob 只看到了 Alice 生活中的 2 年 (参见 Fig.~\ref{fig:twin-paradox-light-signals}). 所以他看到 Alice 的动作变慢, 速率为正常的三分之一. 在返回的 6 年中, Bob 会看到 Alice 在两人重逢前历经的其余 18 年, 因此会看到她的动作加快, 速率为正常的 3 倍.

另一方面, 事件 $(10,8)$ 发出的光直到 $t=18$ 才抵达 Alice (参见 Fig.~\ref{fig:twin-paradox-light-signals}). 所以 Alice 用 18 年看到 Bob 前往行星 X 的 6 年行程, 因而看到他的动作变慢, 速率为正常的三分之一. 在余下的 2 年中, Alice 将看到 Bob 返回的 6 年, 因而看到他的动作加快, 速率为正常的 3 倍.
\RGfigure[.65\textwidth]{twin-paradox-light-signals}{双生子佯谬}
\end{enumerate}

\RGsolution{ex:faster-than-light-paradox}{faster-than-light-paradox}
\begin{enumerate}[label=(\alph*),ref=\alph*,leftmargin=*]
\item 按船员的测量, \emph{Enterprise} 的航程历时
\[
  |(13,12)|=\sqrt{13^2-12^2}=\sqrt{25}=5\ \text{年}.
\]
\item \emph{Enterprise} 的参考系相对地球以速度 $v=\frac{12}{13}$ 运动. 因此 $\sqrt{1-v^2}=\frac5{13}$, 从而
\[
  \begin{cases}
    t'=(1-v^2)^{-1/2}t-v(1-v^2)^{-1/2}x=\frac{13}{5}t-\frac{12}{5}x,\\[2pt]
    x'=(1-v^2)^{-1/2}x-v(1-v^2)^{-1/2}t=\frac{13}{5}x-\frac{12}{5}t.
  \end{cases}
\]
所以在 \emph{Enterprise} 的参考系中, 无线电信号在事件 $(0,0)$ 发出, 导弹在事件 $(28.6,-26.4)$ 发射, 地球在事件 $(2.4,2.4)$ 被摧毁, 而 \emph{Enterprise} 在事件 $(5,0)$ 抵达地球的废墟 (正如应当发生的那样).
\item Fig.~\ref{fig:enterprise-frame-sequence-of-events} 画出了这些事件在 \emph{Enterprise} 参考系中的位置. 事件的先后顺序超乎现实: 地球在 $t'=2.4$ 无缘无故地爆炸; 超光速导弹完好地从爆炸中跃出, 朝 Klingon 行星倒飞, 此时那里正在制造一枚一模一样的复制品; 两枚导弹于 $t'=28.6$ 同时消失, 这一事件本应是地球毁灭的原因, 却发生在 \emph{Enterprise} 已抵达地球废墟之后很久. 这说明, 若允许超光速, 可能得到怎样的荒谬结果.
\RGfigure[.57\textwidth]{enterprise-frame-sequence-of-events}{\emph{Enterprise} 参考系中的事件顺序}
\end{enumerate}
\end{description}
